% Part: normal-modal-logic
% Chapter: completeness
% Section: introduction

\documentclass[../../../include/open-logic-section]{subfiles}

\begin{document}

\olfileid{nml}{com}{int}

\olsection{引言}

如果 $\Sigma$ 是一个模态系统，那么可靠性定理表明：若 $\Sigma \Proves !A$，则 $!A$ 在任何满足如下条件的模型类 $\mClass{C}$ 中都是有效的：$\Sigma$ 中所有公式的所有实例在 $\mClass{C}$ 中都有效。
特别地，这意味着：若 $\Log{K} \Proves !A$，则 $!A$ 在所有模型中为真；若 $\Log{KT} \Proves !A$，则 $!A$ 在所有自反模型中为真；若 $\Log{KD} \Proves !A$，则 $!A$ 在所有持续模型中为真；等等。

完备性是可靠性的逆：例如，\Log{K} 是完备的意味着，如果一个公式~$!A$ 是有效的，则 $\Proves !A$。证明完备性比证明可靠性要困难得多。首先，考虑其逆否命题是有用的：\Log{K} 是完备的当且仅当只要 $\Proves/ !A$，就存在一个反模型，即一个模型~$\mModel{M}$ 使得 $\mSat/{M}{!A}$。等价地（取 $!A$ 的否定），我们可以证明，只要 $\Proves/ \lnot !A$，就存在一个~$!A$ 的模型。在构造这样一个模型时，我们可以利用 $!A$ 所包含的信息。当我们为特定公式寻找模型时，也经常这样做：例如，如果我们想为 $p \lif \Box q$ 找一个反模型，我们知道它必须包含一个 $p$ 为真而 $\Box q$ 为假的世界。而一个 $\Box q$ 为假的世界，意味着必须存在一个从它可达的世界，在其中 $q$ 为假。这就是我们需要知道的全部信息：哪些世界使得命题变元为真，以及哪些世界可从哪些世界可达。

然而，在证明完备性的情况下，我们并没有一个为其构造模型的特定公式~$!A$。我们希望确立的是，对每个满足 $\Proves/[\Sigma] \lnot !A$ 的 $!A$，都存在一个模型。这是一个最低要求，因为\emph{如果} $\Proves[\Sigma] \lnot !A$，根据可靠性，就不存在~$!A$ 的（使得 $\Sigma$ 为真的）模型。现在注意到，$\Proves/[\Sigma] \lnot !A$ 当且仅当 $!A$ 是 $\Sigma$-一致的。（回想一下，$\Sigma \Proves/[\Sigma] \lnot !A$ 与 $!A \Proves/[\Sigma] \lfalse$ 是等价的。）所以我们的任务是为每个 $\Sigma$-一致的公式构造一个模型。

我们将使用的技巧是找到一个 $\Sigma$-一致的公式集合，它包含~$!A$，同时也包含其他一些公式，这些公式告诉我们使 $!A$ 为真的世界必须具有怎样的特征。这样的集合就是\emph{完全的} $\Sigma$-一致集合。仅构造一个只有单个世界的模型来使~$!A$ 为真是不够的，它必须包含多个世界和一个可达关系。包含~$!A$ 的完全 $\Sigma$-一致集合也会包含形如 $\Box !B$ 和~$\Diamond !C$ 的其他公式。在所有可达世界中，$!B$ 必须为真；在至少一个可达世界中，$!C$ 必须为真。为了实现这一点，我们只需简单地取\emph{所有}可能的完全 $\Sigma$-一致集合作为世界集合的基础。棘手之处在于确定在我们的模型中，何时一个完全的 $\Sigma$-一致集合应当被视为可从另一个到达。

我们将证明，在如此定义的模型中，$!A$ 在某个世界（该世界也是一个完全的 $\Sigma$-一致集合）中为真，当且仅当 $!A$ 属于该集合。如果 $!A$ 是 $\Sigma$-一致的，它将属于至少一个完全的 $\Sigma$-一致集合（我们将证明这一事实），因此就会存在一个世界使得 $!A$~为真。这样，我们将得到一个单一的模型，其中每个 $\Sigma$-一致的公式~$!A$ 都在某个世界为真。这个单一的模型就是~$\Sigma$ 的\emph{典范}模型。

\end{document}